\section{Understanding Weaknesses in CNs}

Under our threat model, attackers can reach core network components from outside the trusted boundary, as shown in prior work~\cite{zhang2025invade}, invalidating the traditional isolation assumption that internal components only receive well-formed, benign messages. As a result, we discovered that implicit-trust-driven minimal validation becomes an exploitable attack surface. We made this observation through an empirical analysis of security issues reported in CN implementations’ GitHub repositories. Specifically, we extracted the root cause of each issue, clustered issues with similar causes, and distilled recurring vulnerability patterns. We then utilized these patterns to guide an LLM in identifying previously unseen iTrue vulnerabilities across the codebases.

\noindent\textbf{GitHub security issues in CN}. 
To study the relations among CN-related known vulnerabilities, we curated a dataset of previously reported GitHub security issues in CN. Specifically, we gathered 22 security issues from two widely used open-source CN implementations, Open5GS and free5GC. These reports follow a largely consistent template, typically including a title and bug description, the affected version and components, steps to reproduce, and supporting artifacts such as logs, expected and observed behavior.
These fields provide context to localize the faulty checks, enabling systematic root-cause analysis across issues.

\vspace{2pt}\noindent\textbf{Root-cause analysis and clustering}. Over the reported security issues, we performed a root-cause analysis using Claude Code to extract three elements from each issue: the message(s) that introduce security risk (the \textit{trigger}), the program location in CN implementations where the risk manifests (the \textit{failure point}), and the underlying vulnerability (the \textit{flaw}). We then summarized these elements and their relations into a \textit{root-cause statement}. For example, in \texttt{ogs\_pfcp\_handle\_create\_urr()}, an out-of-bounds read (failure point) can be triggered by a Session Establishment Request that omits mandatory fields (e.g., Downlink Volume) in the Volume Threshold IE (trigger), exposing a missing mandatory-field check (flaw) and causing a crash.
Although these issues span different messages and code paths, many root-cause statements share the same structure: the implementation consumes untrusted message fields without validating mandatory IEs, lengths, or bounds. For instance, in \texttt{ogs\_pfcp\_extract\_node\_id()}, the code copies Node ID data into a fixed-size buffer using a length taken directly from an unverified IE length field in an Association Setup Request, due to missing IE-length validation. While the specific IE types and messages differ, both this case and the above example stem from the same underlying cause, that is, failure to verify the format and completeness of incoming messages. Guided by these similarities, we leveraged an LLM to infer relations among issues and obtained three distinct clusters. 


\vspace{2pt}\noindent\textbf{Analysis and findings}.
Looking into these clusters, we found that CN components often forgo due diligence in validating message format, message semantics, and resource availability, and instead tend to \textit{blindly} act on messages received from internal peers: 

\vspace{2pt}\noindent$\bullet$ \textit{Missing syntactic validation}: 12 issues involve malformed messages being accepted without checking the presence of mandatory fields or enforcing basic format and length constraints. For example, on receiving a Session Establishment Request containing a Volume Threshold IE, the UPF \textit{blindly} proceeds to copy volume fields based on presence flags, without verifying that all required fields are included. As a result, one can craft a truncated Volume Threshold IE to trigger an out-of-bounds read and crash the UPF.

\vspace{2pt}\noindent$\bullet$ \textit{Missing semantic validation}: 8 issues involve messages being accepted without validating semantic constraints or enforcing state-dependent invariants. For example, upon receiving a Session Establishment Request containing an F-TEID value of \texttt{0}, the UPF \textit{blindly} uses the F-TEID to look up the corresponding PDR without checking whether the value falls within a valid range. As a result, an attacker can craft a PFCP message with an out-of-range F-TEID to trigger an assertion failure and crash the UPF.

\vspace{2pt}\noindent$\bullet$ \textit{Missing checking on resource availability}: 2 issues involve risky resource-handling operations where the recipient allocates requested resources without checking availability. For example, upon receiving a PFCP Session Establishment Request, the UPF \textit{blindly} allocates buffers from its internal pool without verifying that sufficient resources remain. As a result, an attacker can repeatedly send PFCP messages to drain the buffer pool, ultimately triggering a segmentation fault and crashing the UPF.

Across these clusters, the common denominator is the \textit{implicit trust} that a CN component places in messages from its internal peers, rooted in the assumption that the CN is well protected and walled off from untrusted environments. This assumption becomes increasingly fragile as CNs move to the cloud: as discussed earlier, new attack surfaces emerge through exposed CN components and loosely protected tunnel interfaces (Section~\ref{sec: threat model}). 
Consequently, we refer to the risks arising from this legacy ``blind faith'' as \textit{implicit trust errors} (\textit{iTrue}).
This observation inspired us to summarize the defining patterns of iTrues to guide the discovery of previously unknown flaws in CN implementations. Finally, although these known issues are largely memory-safety bugs, in part because they were most likely detected by fuzzing, \textit{we show later that the same patterns enable the discovery of more subtle and more severe vulnerabilities, including those that enable session hijacking} (Section~\ref{sec: case studies}).



